% généré par scripts/06_textes.py
\dosmeta{Année de classement : 382 (directe : lecture partielle d’un nom consulaire). Références : ICUR I 3127 = ILCV 2148. Rattachement selon l’EDB : \emph{Incertae originis} (le volume I des ICVR réunit les inscriptions dont la provenance est perdue). Support : Support non précisé par l’EDB. Fiche EDB : Filippo Piazzolla, première version 2015. Notice du catalogue p.~\pageref{not:232}.}
\txtbloc{Texte transmis (EDCS, état 2022)}{\lnum{1}[- - -]na in dom(ino ?) cultrix\par
\lnum{2}[- - -]P nutrivit\par
\lnum{3}[- - - s]uis\par
\lnum{4}[- - - amatrix paup]e\{ro\}rum\par
\lnum{5}[quae pro tanta] merita accepit\par
\lnum{6}[sepulcrum intra l]imina san(c)torum\par
\lnum{7}[- - -]t accepit\par
\lnum{8}[quod multi cupiu]n(t) et rari accipiun(t)\par
\lnum{9}[- - - Antonio et S]⸢y⸣agrio c(on)s(ulibus)}
\txtbloc{Texte de l’EDB (EDB33081)}{\lnum{1}[- - -]na in dom(ino?) cultrix\par
\lnum{2}[- - -]p nutrivit\par
\lnum{3}[- - - s]uis\par
\lnum{4}[- - -amatrix pau]perorum\par
\lnum{5}[- - - quae pro tanta] merita accepit\par
\lnum{6}[sepulcrum intra l]imina santorum\par
\lnum{7}[- - -]t accepit\par
\lnum{8}[quod multi cupiu]n(t) et rari accipiun(t)\par
\lnum{9}[- - - Antonio et S]uagrio c(on)s(ulibus)}
